\documentclass{situationist-preprint}

\title{Catharsis Reconsidered}
\msid{AES-06-b}
\msfamily{AES-06}
\msdiscipline{Aesthetics}
\mstype{Companion manuscript}
\mskeywords{catharsis; Aristotle; tragedy; purgation; clarification; understanding}

\begin{document}
\maketitle

\begin{abstract}
Aristotle says that tragedy, through pity and fear, accomplishes the catharsis of such emotions. AES-06-a argues that tragedy offers a controlled engagement with loss and that its pleasure is a form of understanding. We ask which reading of catharsis, if any, the account supports. We distinguish the purgative, the purificatory and the clarificatory readings, and argue that only the third is consistent with two features that AES-06-a establishes: that the pleasure of tragedy is reflective rather than immediate, and that it is diminished when a tragedy is felt to be gratuitous. The purgative reading, the most influential since Bernays, predicts neither.
\end{abstract}

\section{Three readings}

The definition of tragedy in the \emph{Poetics} ends with the clause that tragedy, through pity and fear, accomplishes the catharsis of such emotions \citep[1449b24--28]{Aristotle}. Aristotle does not explain the term in the \emph{Poetics}, and the interpretive tradition divides into three readings \citep{Halliwell1986}.

\begin{definition}\label{d:readings}
On the \emph{purgative} reading, catharsis is the discharge of an excess of pity and fear, on the model of a medical purge, with relief as its effect \citep{Bernays1857}. On the \emph{purificatory} reading, it is the refinement of the emotions themselves, so that they come to be felt toward the right objects in the right measure. On the \emph{clarificatory} reading, it is an intellectual clarification of the events represented, achieved through the emotions they arouse \citep{Golden1962}.
\end{definition}

\section{Which reading fits}

Theorem~4.2 of AES-06-a states that the pleasure of tragedy is the pleasure of understanding a loss, achieved through painful emotions whose bounds make them tolerable; it rests on the control principle P1 and on the principle P2 that some understanding of loss is available only through an emotional response to it.

\begin{proposition}\label{p:fit}
Of the three readings, only the clarificatory reading identifies catharsis with the source of the pleasure that Theorem~4.2 of AES-06-a describes.
\end{proposition}

\begin{proof}
On the purgative reading, the effect of tragedy is relief, which is a pleasure in the cessation of a pain. It is immediate rather than reflective, and it would be as great for a gratuitous tragedy as for a profound one, provided the emotions were equally strong; both predictions conflict with the features of tragic pleasure AES-06-a identifies. On the purificatory reading, the effect is a change in the dispositions of the spectator, which may be real but is not a pleasure taken in the work, and it is realized over time rather than in the experience of the tragedy. On the clarificatory reading, the effect is understanding achieved through pity and fear, which by Principle~P2 is what the emotions make possible; it is reflective, and it is diminished when the loss is gratuitous, since a gratuitous loss offers nothing to understand. This is the pleasure of Theorem~4.2.
\end{proof}

The clarificatory reading also fits Aristotle's remark that people delight in imitations because learning is most pleasant, and that they enjoy seeing likenesses because in contemplating them they learn and infer what each thing is \citep[1448b4--17]{Aristotle}. \citet{Nussbaum1986} develops a related reading, on which the emotions of tragedy are themselves modes of recognition of what matters in human life.

\begin{objection}[Philology]
The medical sense of \emph{katharsis} was well established in Aristotle's time, and his discussion of music in the \emph{Politics} uses the term in connection with relief. The purgative reading has textual support that the clarificatory one lacks.
\end{objection}

\begin{reply}
The argument here does not decide what Aristotle meant. It shows which reading is consistent with the best account of the pleasure of tragedy. If Aristotle meant purgation, then on the present account he described an effect tragedy may have but not the one that explains why it is sought. \citet{Lear1988} argued, on independent grounds, that none of the standard readings fits the text without strain; the account of AES-06-a gives a reason to prefer clarification that does not depend on settling the philology.
\end{reply}

\section{Conclusion}

If catharsis names the effect in virtue of which tragedy pleases, it is clarification: an understanding of loss achieved through pity and fear. The purgative reading describes an effect that is neither reflective nor sensitive to whether a tragedy's suffering is gratuitous, and so cannot be the source of tragic pleasure.

\begin{thebibliography}{99}\small
\bibitem[Aristotle(Poet.)]{Aristotle} Aristotle. \emph{Poetics}. Cited by Bekker page. Trans. S. Halliwell, Loeb Classical Library, Harvard University Press, 1995.
\bibitem[Bernays(1857)]{Bernays1857} Bernays, J. (1857). \emph{Grundz\"uge der verlorenen Abhandlung des Aristoteles \"uber Wirkung der Trag\"odie}. Trewendt.
\bibitem[Golden(1962)]{Golden1962} Golden, L. (1962). Catharsis. \emph{Transactions and Proceedings of the American Philological Association}, 93, 51--60.
\bibitem[Halliwell(1986)]{Halliwell1986} Halliwell, S. (1986). \emph{Aristotle's Poetics}. Duckworth.
\bibitem[Lear(1988)]{Lear1988} Lear, J. (1988). Katharsis. \emph{Phronesis}, 33(3), 297--326.
\bibitem[Nussbaum(1986)]{Nussbaum1986} Nussbaum, M.~C. (1986). \emph{The Fragility of Goodness: Luck and Ethics in Greek Tragedy and Philosophy}. Cambridge University Press.
\end{thebibliography}

\end{document}
